\documentclass[../../main.text]{subfiles}
%Deutsche Sprache
\usepackage[ngerman]{babel}
\usepackage[T1]{fontenc}
\usepackage{lmodern}
\usepackage[thmmarks,amsmath,hyperref,noconfig]{ntheorem}

%Zeilenabstand
\usepackage[onehalfspacing]{setspace}

\begin{document}
\begin{onehalfspacing}
Im letzten Jahrzehnt ist der schonende Umgang mit Ressourcen wichtiger 
geworden und kann als zentraler Bestandteil der Kosten einer Dokumentenarchivierung 
angesehen werden. Die voranschreitende Digitalisierung 
erfordert immer mehr und immer effektiver genutzten Speicher.
Wie schnell der Papierstapel wächst und dabei die Übersicht und 
die Ordnung beizubehalten, weiß jeder, der sich um Versicherungen, 
Ärzte oder amtliche Angelegenheiten kümmern muss.
Es ist durchaus denkbar, dass ein mittelständisches Unternehmen 
ein Vielfaches dessen pro Stunde verarbeitet. 
Eine manuelle Bearbeitung des Papierstapels wird mit zunehmender 
Größe auch zunehmend unpraktisch, da die Effizienz abnimmt. 
Daher gibt es im Enterprise Content Management-Sektor
viele Produkte, die bei der Speicherung und der Verwaltung 
ebendieser Dokumente effektiv unterstützen. 
Die maximale Speicherkapazität und die schnelle Erreichbarkeit sind 
hierbei von äußerster Bedeutung. Alfresco, OpenText Content Suite oder auch ImageMaster
sind Beispiele für Enterprise Content Management-Lösungen. 
Alle diese Lösungen haben als Ziel eine 
möglichst hohe Leistungsfähigkeit und hohe Skalierbarkeit zu gewährleisten.
Besonders in Unternehmen mit sehr hohen Zugriffszahlen sind die Stabilität und
die Effizienz der Enterprise Content Management-Lösung von allerhöchster Priorität.
Um ein vorhandenes oder ein völlig neues System auf Stabilität
zu testen, werden Last-Tests verwendet. Wenn die maximal erbringbare Kapazität
eines Systems bekannt ist, lässt sich das System ressourcenschonender betreiben
und gewinnbringender verkaufen.
%%%%%%%%%%%%%%%%%%%%%%%%%%%%%%%%%%%%%%%%%%%%%%%%
\end{onehalfspacing}
\end{document}
%%%%%%%%%%%%%%%%%%%%%%%%%%%%%%%%%%%%%%%%%%%%%%%